\documentclass[12pt,a4paper]{article}
\usepackage{cmap} % Makes the PDF copiable. See http://tex.stackexchange.com/a/64198/25761
\usepackage[T1]{fontenc}
\usepackage[brazil]{babel}
\usepackage[utf8]{inputenc}
\usepackage{amsmath}
\usepackage{amsfonts}
\usepackage{amssymb}
\usepackage{amsthm}
\usepackage{textcomp} % \degree
\usepackage{gensymb} % \degree
\usepackage[usenames,svgnames,dvipsnames]{xcolor}
\usepackage{hyperref}
\usepackage{graphicx}
\usepackage[margin=2cm]{geometry}

\hypersetup{
    colorlinks = true,
    allcolors = {blue}
}

% TODO: Consider using exsheets
% http://linorg.usp.br/CTAN/macros/latex/contrib/exsheets/exsheets_en.pdf
%
% http://ctan.org/tex-archive/macros/latex/contrib/exercise/
% Options: answerdelayed,lastexercise,noanswer
\usepackage[answerdelayed,lastexercise]{exercise}

\addto\captionsbrazil{%
\def\listexercisename{Lista de exerc\'icios}%
\def\ExerciseName{Exerc\'icio}%
\def\AnswerName{Solu\c{c}\~ao do exerc\'icio}%
\def\ExerciseListName{Ex.}%
\def\AnswerListName{Solu\c{c}\~ao}%
\def\ExePartName{Parte}%
\def\ArticleOf{de\ }%
}

\renewcommand{\ExerciseHeaderTitle}{(\ExerciseTitle)\ }
\renewcommand{\ExerciseListHeader}{%\ExerciseHeaderDifficulty%
\textbf{%\ExerciseListName\
\ExerciseHeaderNB.\ %
%\ --- \
\ExerciseHeaderTitle}%
%\ExerciseHeaderOrigin
\ignorespaces}
\renewcommand{\AnswerListHeader}{\textbf{\ExerciseHeaderNB.\ (\AnswerListName)\ }}

\newcommand*\sen{\operatorname{sen}}
\newcommand*\cotg{\operatorname{cotg}}
\newcommand*\cosec{\operatorname{cossec}}
\newcommand*\dom[1]{\operatorname{Dom}\left(#1\right)}
\newcommand*\im[1]{\operatorname{Im}\left(#1\right)}

\newcommand*\R{\mathbb{R}}

\renewcommand{\theenumi}{\alph{enumi}}
\renewcommand\labelenumi{(\theenumi) }

\newcommand*\tipo{Prova I}
\newcommand*\turma{TADS121-01U}
\newcommand*\disciplina{CDI0001}
\newcommand*\eu{Helder G. G. de Lima}
\newcommand*\data{25/09/2017}

\author{\eu}
\title{\tipo - \disciplina}
\date{\data}

\begin{document}
\thispagestyle{empty}
\newgeometry{margin=2cm,bottom=0.5cm}
\begin{center}
\includegraphics[width=9.0cm]{marca} \\
\textbf{\tipo\ (\disciplina / \turma)} \\
Prof. \eu\footnote{
Este é um material de acesso livre distribuído sob os termos da licença \href{https://creativecommons.org/licenses/by-sa/4.0/deed.pt_BR}{Creative Commons BY-SA 4.0}}
\end{center}

\noindent Nome do(a) aluno(a): \underline{\hspace{9,7cm}} Data: \underline{\data}

%\section*{Instruções}
\begin{center}\fbox{
\begin{minipage}{14cm}

{\footnotesize
\begin{itemize}
\renewcommand{\theenumi}{\Roman{enumi}}
\item Identifique-se em todas as folhas.
\item Mantenha o celular e os demais equipamentos eletrônicos desligados durante a prova.
\item Justifique cada resposta com cálculos ou argumentos baseados na teoria estudada.
\item Escolha apenas 5 das 7 questões para serem corrigidas (totalizando 10,0 pontos).
\end{itemize}
}

\end{minipage}
}
\end{center}

%\section*{Questões}
\begin{ExerciseList}
\Exercise[title={2,0}] Sejam $f(x) = \sqrt{x}$, $g(x) = 2 - |3 - x|$ e $h(x) = |x|$.
\begin{enumerate}
\item Faça um esboço dos gráficos de $g$ e $h$.
\item Obtenha o domínio de $f \circ g$.
\end{enumerate}
\Answer
\begin{enumerate}
\item Uma vez esboçado o gráfico de $h$, cuja fórmula é mais simples, pode-se construir o de $g$ transladando-o 3 unidades para a direita, refletindo o resultado em relação ao eixo $x$ e transladando duas unidades para cima:
\begin{center}
\includegraphics[width=10cm]{img/prova-1-tads-graph-g-h}
\end{center}

\item Como $(f \circ g) (x) = f(g(x)) = \sqrt{ 2 - |3 - x| }$ e $\sqrt{u}$ só está definida para $u \geq 0$, os únicos valores de $x \in \R$ que estão no domínio de $f \circ g$ são aqueles para os quais $2 - |3 - x| \geq 0$. Tais valores podem ser determinados assim:
\[
2 - |3 - x| \geq 0
\Leftrightarrow
|3 - x| \leq 2
\Leftrightarrow
-2 \leq 3 - x \leq 2
\Leftrightarrow
-5 \leq - x \leq -1
\Leftrightarrow
5 \geq x \geq 1
\]
Portanto, $\dom{ f\circ g } = [1, 5]$.
\end{enumerate}

\Exercise[title={2,0}] Seja $q$ uma função definida pela regra $q(x) = 3 + \dfrac{x^3-1}{x^3+1}$.
\begin{enumerate}
\item Obtenha a função inversa de $q$, deixando explícita a fórmula para $q^{-1}(x)$.
\item Obtenha o domínio e a imagem de $q$ e de $q^{-1}$.
\end{enumerate}
\Answer
\begin{enumerate}
\item Tem-se:
\begin{align*}
y = q(x) = 3 + \dfrac{x^3-1}{x^3+1}
& \Leftrightarrow
y-3 = \dfrac{x^3-1}{x^3+1}
\Leftrightarrow
(y-3)(x^3+1) = x^3-1\\
& \Leftrightarrow
yx^3+y-3x^3-3 = x^3-1
\Leftrightarrow
yx^3-4x^3 = 2-y\\
& \Leftrightarrow
(y-4)x^3= 2-y
\Leftrightarrow
x^3= \frac{2-y}{y-4}
\Leftrightarrow
x= \sqrt[3]{ \frac{2-y}{y-4} }
\end{align*}
Portanto, usando $x$ para a variável independente, como é de costume, tem-se $q^{-1}(x) = \sqrt[3]{ \frac{2-x}{x-4} }$.
\item Com base no item anterior, $q^{-1}(x)$ está bem definida desde que o denominador $x-4$ não seja nulo, isto é, para todo $x \neq 4$. Logo $\im{q} = \dom{q^{-1}} = \R \setminus \{ 4 \}$. Analogamente, $q(x)$ está definida para todo $x$ que não anule o denominador $x^3+1$, isto é, para $x \neq -1$. Assim, $\im{q^{-1}} =\dom{q} = \R \setminus \{ -1 \}$.
\end{enumerate}

\Exercise[title={2,0}] Seja $p(x) = x^2 - 6x$. Determine o valor (ou valores) de $a \in \R$ para que a função definida por $t(x) = p(x + a)$ seja uma função par.
\Answer Uma função $t$ é par caso $t(-x) = t(x), \forall x \in \R$. Então:
\begin{align*}
t(x) = t(-x)
\Leftrightarrow
p(x+a) & = p(-x+a)\\
\Leftrightarrow
(x+a)^2 - 6(x+a) & = (-x+a)^2 - 6(-x+a)\\
\Leftrightarrow
x^2 + 2ax + a^2 - 6x - 6a & = x^2 - 2ax + a^2 + 6x - 6a \\
\Leftrightarrow
2ax - 6x & = - 2ax + 6x
\Leftrightarrow
4ax = 12x
\Leftrightarrow
4a = 12, \text{ para todo } x \neq 0.
\end{align*}
Assim, é preciso que $a = 12/4 = 3$ para que $t$ seja uma função par.

\Exercise[title={2,0}] Utilize a \textit{definição} de limite para verificar que $\displaystyle \lim_{x \to 1} \dfrac{(x+1)(x^2 + x - 2)}{x^2 - 1} = 3$.
\Answer Para que seja verdade que  $\displaystyle \lim_{x \to 1} \dfrac{(x+1)(x^2 + x - 2)}{x^2 - 1} = 3$, é preciso que para cada $\varepsilon > 0$ exista algum $\delta > 0$ tal que $\left|\dfrac{(x+1)(x^2 + x - 2)}{x^2 - 1} - 3\right| < \varepsilon$ sempre que $0 < |x-1| < \delta$. Considerando que
\begin{align*}
\dfrac{(x+1)(x^2 + x - 2)}{x^2 - 1} - 3
& = \dfrac{(x+1)(x^2 + x - 2)}{(x+1)(x-1)} - 3
  = \dfrac{x^2 + x - 2}{x-1} - 3
  = \dfrac{x^2 + x - 2 - 3x + 3}{x-1} \\
& = \dfrac{x^2 -2x +1}{x-1}
  = \dfrac{(x-1)^2}{x-1}
  = x-1,
\end{align*}
conclui-se que é possível obter $|x-1| < \varepsilon$ quando $0 < |x-1| < \delta$ bastando escolher $\delta = \varepsilon$.

\Exercise[title={2,0}] Calcule, se existir, $\displaystyle\lim_{x \to 0} \cosec(x) - \cotg(x)$.
\Answer Tem-se:
\begin{align*}
\lim_{x \to 0} \cosec(x) - \cotg(x)
& = \lim_{x \to 0} \frac{1}{\sen(x)} - \frac{\cos(x)}{\sen(x)}
  = \displaystyle\lim_{x \to 0} \frac{1-\cos(x)}{\sen(x)}
    \cdot
    \frac{(1+\cos(x))}{(1+\cos(x))} \\
& = \lim_{x \to 0} \frac{1-\cos^2(x)}{\sen(x)(1+\cos(x))}
  = \lim_{x \to 0} \frac{\sen^2(x)}{\sen(x)(1+\cos(x))} \\
& = \lim_{x \to 0} \frac{\sen(x)}{1+\cos(x)}
  = \lim_{x \to 0} \frac{0}{1+1} = 0
\end{align*}

\Exercise[title={2,0}] Calcule, se existir,
$\displaystyle\lim_{x \to +\infty}
       \frac{(e^{1/x})^2-1}
            {e^{1/x}-1}$.
\Answer Tem-se:
\begin{align*}
  \lim_{x \to +\infty} \frac{(e^{1/x})^2-1}{e^{1/x}-1}
= \lim_{u \to 0^+} \frac{(e^u)^2-1}{e^u-1}
= \lim_{v \to 1} \frac{v^2-1}{v-1}
= \lim_{v \to 1} \frac{(v-1)(v+1)}{v-1}
= \lim_{v \to 1} v+1
= 2
\end{align*}
\Exercise[title={2,0}] Verifique se há algum $k\in \R$ para o qual a função a seguir seja continua em $x = 4$:
\[
f(x) =
\begin{cases}
\dfrac{x - 4}{x^2 - 4x}, & \text{se } x < 4 \\
k, & \text{se } x = 4 \\
\dfrac{\sqrt{x} - 2}{x - 4}, & \text{se } x > 4
\end{cases}
\]
\Answer Para que $f$ seja contínua em $x = 4$, é preciso que exista $\lim_{x \to 4} f(x)$ e que ele coincida com $f(4) = k$. A existência do limite estará garantida se os limites laterais forem iguais:
\[
\lim_{x \to 4^-} \dfrac{x - 4}{x^2 - 4x}
= \lim_{x \to 4^-} \dfrac{x - 4}{x(x - 4)}
= \lim_{x \to 4^-} \dfrac{1}{x}
= 1 / 4
\]
e
\[
\lim_{x \to 4^+} \dfrac{\sqrt{x} - 2}{x - 4}
= \lim_{x \to 4^+} \dfrac{(\sqrt{x} - 2)(\sqrt{x} + 2)}{(x - 4)(\sqrt{x} + 2)}
= \lim_{x \to 4^+} \dfrac{x - 4}{(x - 4)(\sqrt{x} + 2)}
= \lim_{x \to 4^+} \dfrac{1}{\sqrt{x} + 2}
= 1 / 4
\]
Logo, $f$ será contínua se $k = 1/4$.
\end{ExerciseList}

\begin{center}
BOA PROVA!
\end{center}

\newpage
\restoregeometry
\section*{Respostas}
\shipoutAnswer
\end{document}
